% 来源：标题摘要关键词_Loop_V3.docx。
\begin{center}
{\zihao{3}\heiti\bfseries 中药材热风烘干过程的传热传质建模与烘干时长预测}
\end{center}
\vspace{0.5em}
\begin{center}{\heiti\zihao{4} 摘\quad 要}\end{center}

本文通过有限体积法（FVM）进行空间离散，并采用变阶、自适应步长的后向差分公式（BDF）进行时间积分，对圆柱形中药材的热风烘干过程进行分析。逐步考虑密度、比热容、导热系数和水分扩散系数的变化及径向收缩，着重分析中部截面的温度和含水率分布。

\par\vspace{\baselineskip}
{\heiti\bfseries 针对问题一，}在药材尺寸及密度、比热容、导热系数保持不变、水分扩散系数随含水率变化的条件下，分别建立温度场和含水率场模型。使用分段三次Hermite插值（PCHIP）\cite{fritsch1980monotone}，得到环境温度和水分浓度关于时间的连续函数，并采用对流换热和对流传质边界，求解预热阶段的场量分布。数值结果表明，在$1800\,\mathrm{s}$时，中部截面中心与表面温度分别为$33.5758\unitC$和$36.7858\unitC$，含水率分别为$2.5500\,\mathrm{kg/kg}$和$1.5103\,\mathrm{kg/kg}$。

\par\vspace{\baselineskip}
{\heiti\bfseries 针对问题二，}在药材尺寸保持不变的条件下，根据当前温度和含水率更新密度、比热容、导热系数和水分扩散系数，联立求解温度场和含水率场。数值结果表明，在$3\,\mathrm{h}$时，中心与表面温度分别为$49.8506\unitC$和$49.9674\unitC$，仅相差$0.1168\unitC$；含水率分别为$1.7662\,\mathrm{kg/kg}$和$1.0081\,\mathrm{kg/kg}$，相差$0.7581\,\mathrm{kg/kg}$。药材的热平衡明显早于水分平衡，后续烘干主要受内部水分向外扩散控制。

\par\vspace{\baselineskip}
{\heiti\bfseries 针对问题三，}将附件环境数据末20组成对观测的平均值作为长时烘干的环境条件，以全域最大含水率首次降至$0.15\,\mathrm{kg/kg}$的时刻作为干燥终点。计算得到烘干总时间为$57.4354\,\mathrm{h}$；全域扫描结果显示，几何中心最后达到干燥阈值。

\par\vspace{\baselineskip}
{\heiti\bfseries 针对问题四，}在径向收缩的动态区域上重新求解温度场和含水率场，并沿用问题三的环境延拓条件和干燥判据。计算得到烘干时间为$51.0583\,\mathrm{h}$，结束时药材半径为$1.2000\,\mathrm{cm}$。与问题三相比，烘干时长减少$6.3772\,\mathrm{h}$，即$11.10\%$；该时间差反映了半径变化及密度、比热容、导热系数和水分扩散系数变化的共同作用。

\par\vspace{\baselineskip}
\noindent\textbf{关键词：}中药材烘干；传热；传质；有限体积法；径向收缩
